\documentclass[sigplan,screen,nonacm]{acmart}

\usepackage{graphicx}
\usepackage{float}
\usepackage{booktabs}
\usepackage{amsmath}
\usepackage{microtype}
\usepackage{xcolor}

\setcopyright{none}
\settopmatter{printacmref=false, printfolios=false}

\title{How Much Text Is Enough?\\
Measuring the Minimum Amount of Text Required\\for LLM-Based Deanonymization}

\author{Gleb Petrov}
\affiliation{\institution{University of Luxembourg}}

\begin{document}

\begin{abstract}
Large Language Models (LLMs) can link pseudonymous online users from writing alone,
yet the quantitative relationship between text volume and attack success remains
underexplored. We study this question through a controlled truncation experiment on
Reddit (POOL-EN) and Hacker News using the Extract--Search--Reason--Calibrate (ESRC)
pipeline across eight truncation levels (T1--T8).

Our most striking finding is a severe \textbf{confidence miscalibration} in GPT-4o:
despite top-1 accuracy approaching 40\%, \textbf{Recall@90\% Precision remains near
zero} on both platforms (Reddit $\leq$2.1\%; Hacker News $\leq$0.62\%). At a stated
confidence of 0.85, the model is empirically correct only 26.7\% of the
time---making confidence-gated deployment largely ineffective.

Top-1 accuracy nonetheless increases nearly monotonically with text volume, from
8.8\% to 39.4\% on Reddit and from 4.8\% to 38.0\% on Hacker News (Spearman
$\rho = 0.976$, $p < 0.001$ on Reddit; $\rho = 1.000$, $p < 0.001$ on Hacker News,
computed over truncation level ranks T1--T8). Bootstrap 95\% confidence intervals overlap
substantially across adjacent levels, indicating that small stepwise differences
should not be overinterpreted. Exploratory changepoint analysis (PELT, RBF) suggests a
possible transition near 50 comments under one penalty setting, but this estimate is
sensitive to the penalty parameter and should be treated as illustrative rather than
definitive. Finally, conditional Reason-stage success given retrieval of the true
match in the top-15 candidates reaches 64--73\% at T8, identifying reasoning---not
search---as the primary pipeline bottleneck.

Two additional experiments at fixed volume (T4, 50 comments) show that \textit{what}
users write appears to matter more than \textit{where} they post. Classifying full query
histories into Personal, Opinion, and Topical comments, Personal-only profiles
achieve 32.4\% top-1 accuracy versus 18.8\% (Opinion) and 18.3\% (Topical;
$p < 0.00001$ for Personal vs.\ both). A community-diversity grouping by unique
subreddits shows a non-significant inverse trend (low-diversity users appear more
identifiable than high-diversity users), but bootstrap intervals overlap and pairwise
tests do not reach significance at $\alpha = 0.05$.
\end{abstract}

\maketitle


\section{Introduction}


Online platforms accumulate vast quantities of user-generated text. Even when
explicit identifiers are removed, the content itself---topic preferences, writing
style, behavioral patterns, and idiosyncratic vocabulary---may be sufficient to
re-identify a user across different contexts or time periods.

Recent work by Lermen et al.~\cite{lermen2026} demonstrated that LLMs can be used
to link pseudonymous users at scale through an automated Extract--Search--Reason
pipeline. Their results show that deanonymization is not only possible but
surprisingly effective even within large candidate pools. However, a fundamental
question was left open: \textit{how much text must a user produce before they
become identifiable?}

This question matters for practical privacy reasoning. A user who posts five
comments occupies a very different risk profile than one with five hundred. Understanding
where the threshold lies---and whether it differs across platforms---provides
actionable guidance for users, platform designers, and policymakers.

De~Montjoye et al.~\cite{demontjoye2013} showed that a small number of mobility
records can make an individual \emph{uniquely identifiable} in a dataset---an analogy
we borrow only to motivate the question of \emph{minimum sufficient identifying
information}. Our work asks the textual equivalent: how many comments constitute
enough evidence for LLM-based re-identification in a closed candidate pool? We do
not equate dataset uniqueness with deanonymization accuracy.

We make the following contributions:
\begin{itemize}
  \item A systematic truncation study across eight text volume levels on two
        independent platforms (Reddit and Hacker News), with exploratory
        changepoint-based estimates of \textbf{Minimum Sufficient Text (MST)} and
        penalty-sensitivity analysis.
  \item An analysis of \textbf{confidence miscalibration} in GPT-4o, showing that
        stated confidence substantially overstates empirical accuracy, with direct
        implications for the practical feasibility of LLM-based attacks.
  \item A decomposition of pipeline performance into \textbf{Search-stage} and
        \textbf{Reason-stage} contributions, revealing that retrieval is not the
        primary bottleneck.
  \item A \textbf{content-type} experiment showing that personal disclosures are
        significantly more identifying than opinion or topical comments at fixed
        volume (32.4\% vs.\ $\sim$18\% top-1 accuracy; $p < 0.00001$).
  \item A \textbf{community-diversity} experiment revealing a non-significant inverse
        trend in which users active in fewer subreddits appear more deanonymizable,
        with overlapping bootstrap confidence intervals.
\end{itemize}


\section{Related Work}


\paragraph{LLM-based deanonymization.}
Lermen et al.~\cite{lermen2026} introduced the ESRC pipeline and demonstrated large-scale de-anonymization on Reddit. Their Figure~4b shows that recall increases
with the number of movies a user discusses---a single observation that motivates
our systematic study. Our work generalizes this finding across text volumes,
platforms, and evaluation metrics.

\paragraph{Minimum identifying information.}
De~Montjoye et al.~\cite{demontjoye2013} established that human mobility traces are
highly unique at small sample sizes. We adopt only the conceptual framing---how
little data suffices before re-identification risk rises materially---and
operationalize a textual minimum sufficient set via truncation and changepoint
analysis, without mapping uniqueness rates to attack success rates.

\paragraph{$k$-Anonymity and privacy models.}
Sweeney~\cite{sweeney2002} formalized the notion of $k$-anonymity---the idea that
a record is private only if it is indistinguishable from at least $k-1$ others.
Our empirical results complement this theoretical framework by quantifying how
much text is needed to break anonymity in practice.


\section{Methodology}


We adopt the ESRC framework of Lermen et al.~\cite{lermen2026} without modification,
using their published prompts verbatim to ensure comparability.

\subsection{Extract}
Raw user comments are summarized into compact trait-based profiles using GPT-4o-mini
with the Lermen G.2 summarization prompt. The resulting summaries capture recurring
interests, writing style, demographic indicators, and behavioral characteristics.

\subsection{Search}
Summaries are embedded using \texttt{all-mpnet-base-v2} (SentenceTransformers) and
indexed with FAISS. For each query profile, the top-15 most similar candidate profiles
are retrieved by cosine similarity.

\subsection{Reason}
GPT-4o receives the query summary alongside the 15 retrieved candidates and selects
the most likely match, returning a confidence score in $[0, 1]$.

\subsection{Calibrate}
Predictions are ranked by confidence and evaluated using Precision-Recall curves.
Our primary metrics are:
\begin{itemize}
  \item \textbf{Recall@90\% Precision}: maximum recall while maintaining $\geq 90\%$ precision.
  \item \textbf{Recall@99\% Precision}: same, at stricter precision.
  \item \textbf{Top-1 Accuracy}: fraction of queries correctly matched (no threshold).
\end{itemize}

\subsection{Truncation Levels}
For each query user, we construct eight versions of the query profile by retaining
only the first $N$ comments in chronological order, simulating an attacker with
access only to a user's early activity:

\begin{table}[H]
\centering
\caption{Truncation levels.}
\begin{tabular}{ccc}
\toprule
Level & Comments & Approx.\ word count \\
\midrule
T1 & 5   & 200--300 \\
T2 & 10  & 400--600 \\
T3 & 25  & 1{,}000--1{,}500 \\
T4 & 50  & 2{,}000--3{,}000 \\
T5 & 100 & 4{,}000--6{,}000 \\
T6 & 200 & 8{,}000--12{,}000 \\
T7 & 500 & 20{,}000--30{,}000 \\
T8 & Full history & --- \\
\bottomrule
\end{tabular}
\end{table}

\subsection{Changepoint Detection}
We apply the PELT algorithm~\cite{truong2020} with an RBF cost function (via the
\texttt{ruptures} library) to the sequence of top-1 accuracy values across T1--T8.
The detected changepoint defines the \textbf{Minimum Sufficient Text (MST)}:
the truncation level beyond which marginal gains in deanonymization performance
diminish substantially.

\subsection{Diversity Experiment}
To test whether community breadth modulates risk independently of text volume, we
group Reddit query users by the number of unique subreddits represented in their T4
(50-comment) profiles: \textbf{low} diversity (1--3 subreddits; $n = 58$),
\textbf{medium} (4--10; $n = 179$), and \textbf{high} (11+; $n = 263$). Each group
uses the same T4 comment count; only the breadth of communities differs. We reuse
candidate summaries and embeddings from the main Reddit ESRC pipeline and run the
full ESRC procedure per group with GPT-4o-mini for Extract and GPT-4o for Reason.

\subsection{Content-Type Experiment}
We classify each user's full 500-comment query history (not T4 alone) into three
content types using GPT-4o-mini: \textbf{Personal (P)}---disclosure of facts about
the author's life; \textbf{Opinion (O)}---beliefs, preferences, and value
judgments; and \textbf{Topical (T)}---topic discussion without personal disclosure.
Requiring $\geq 50$ comments of each type is impossible from only 50 T4 comments, so
classification uses the full chronology. For each qualifying user, we construct
type-filtered profiles from the first 50 comments of each type. Candidate summaries
and embeddings are reused from the main Reddit pipeline; Reason uses GPT-4o. The ESRC
cohort comprises users with successful Extract summaries for all three types
($n = 432$ of 447 users qualifying on all three types; 15 Extract calls failed on
oversized profiles).


\section{Experimental Setup}


\subsection{Reddit (POOL-EN)}
\begin{itemize}
  \item 500 query users, 1{,}000 candidate users
  \item 500 matched pairs + 500 distractor candidates
  \item Users filtered to $\geq 500$ comments per temporal half, $\geq 4$ years
        of activity, excluding bots and deleted accounts
  \item Temporal split: 365-day gap between query and candidate halves
\end{itemize}

\subsection{Hacker News}
\begin{itemize}
  \item 500 query users, 1{,}000 candidate users
  \item Same filtering and splitting criteria as Reddit
  \item Data accessed via Google BigQuery public dataset
\end{itemize}

Both datasets use the same ESRC pipeline, prompts, and evaluation procedure,
enabling direct cross-platform comparison.


\section{Results}


\subsection{Top-1 Accuracy Across Truncation Levels}

\begin{table}[H]
\centering
\caption{Top-1 deanonymization accuracy across truncation levels.}
\begin{tabular}{cccc}
\toprule
Level & Comments & Reddit (\%) & Hacker News (\%) \\
\midrule
T1 & 5   & 8.8  & 4.8  \\
T2 & 10  & 13.4 & 8.6  \\
T3 & 25  & 16.6 & 15.6 \\
T4 & 50  & 25.2 & 22.2 \\
T5 & 100 & 28.2 & 26.0 \\
T6 & 200 & 33.2 & 32.4 \\
T7 & 500 & 39.8 & 34.6 \\
T8 & Full & 39.4 & 38.0 \\
\bottomrule
\end{tabular}
\end{table}

Performance increases nearly monotonically on both platforms. Reddit consistently
outperforms Hacker News at lower truncation levels, though bootstrap 95\% confidence
intervals overlap at most adjacent levels (Table~\ref{tab:bootstrap_ci}), so modest
stepwise gaps should not be treated as definitive. The platforms converge near
38--40\% at full profile size. The difference between T7 and T8 falls within the
expected sampling uncertainty and should not be interpreted as evidence of a genuine
performance decrease. Overall, T8 top-1 accuracy is approximately 4.5$\times$ that
at T1 on Reddit (39.4\%/8.8\%) and 7.9$\times$ on Hacker News (38.0\%/4.8\%).

\begin{table}[H]
\centering
\caption{Bootstrap 95\% confidence intervals for top-1 accuracy (10{,}000 resamples).}
\label{tab:bootstrap_ci}
\begin{tabular}{lccc}
\toprule
Level & Reddit (\%) & Hacker News (\%) \\
\midrule
T1 & 8.8 [6.4, 11.4]   & 4.8 [3.0, 6.8]   \\
T4 & 25.2 [21.4, 29.0] & 22.2 [18.8, 25.8] \\
T7 & 39.8 [35.6, 44.3] & 34.6 [30.4, 38.8] \\
T8 & 39.4 [35.2, 43.6] & 38.0 [33.8, 42.2] \\
\bottomrule
\end{tabular}
\end{table}

\begin{figure}[H]
\centering
\includegraphics[width=0.85\linewidth]{reddit_vs_hn_accuracy_with_ci.png}
\caption{Top-1 deanonymization accuracy across truncation levels with bootstrap 95\%
confidence intervals (10{,}000 resamples). Shaded bands indicate substantial
uncertainty; adjacent-level differences are often not statistically separable at
this sample size.}
\end{figure}

\subsection{Spearman Correlation}

The monotonic relationship between truncation level and accuracy is statistically robust
when truncation level ranks are used as the independent variable (T1--T8 coded as
$1, 2, \ldots, 8$):

\begin{itemize}
  \item Reddit: $\rho = 0.976$, $p < 0.001$
  \item Hacker News: $\rho = 1.000$, $p < 0.001$
\end{itemize}

These values confirm a strong rank correlation between truncation level and
deanonymization performance across both platforms.

\subsection{Marginal Efficiency of Additional Text}

To understand \textit{where} additional text provides the most value to an attacker,
we compute the accuracy gain per additional comment at each truncation step:

\begin{table}[H]
\centering
\caption{Marginal accuracy gain per additional comment (Reddit).}
\begin{tabular}{cccc}
\toprule
Transition & Comments & $\Delta$ Accuracy & Comment gain \\
\midrule
T1$\to$T2 & 5   & +4.6\% & 0.92\%/comment \\
T2$\to$T3 & 15  & +3.2\% & 0.21\%/comment \\
T3$\to$T4 & 25  & +8.6\% & 0.34\%/comment \\
T4$\to$T5 & 50  & +3.0\% & 0.06\%/comment \\
T5$\to$T6 & 100 & +5.0\% & 0.05\%/comment \\
T6$\to$T7 & 300 & +6.6\% & 0.02\%/comment \\
T7$\to$T8 & --- & $-$0.4\% & --- \\
\bottomrule
\end{tabular}
\end{table}

The marginal return on additional text declines sharply after T2 and stabilizes at
near-zero beyond T5. From the attacker's perspective, the first 10 comments are by
far the most information-dense, yielding 0.92\% accuracy gain per comment---approximately
2.7 times the gain rate of the T3--T4 transition (0.92/0.34).

\subsection{Minimum Sufficient Text (MST)}

We treat MST as an \emph{exploratory} estimate rather than a discovered constant.
Applying PELT changepoint detection (RBF cost) to the top-1 accuracy sequence yields
breakpoint sets that depend strongly on the penalty parameter (Table~\ref{tab:mst_pen}). At
penalty~$= 2$, both platforms show a single breakpoint at T4 (50 comments); at
penalty~$= 0.5$--$1$, multiple breakpoints appear (e.g., T4 and T6 on Reddit); at
penalty~$\geq 3$, no breakpoints are detected. One plausible reading of the data is
therefore an early acceleration near 50 comments followed by diminishing returns,
but this interpretation is not stable across penalty choices.

\begin{table}[H]
\centering
\caption{PELT breakpoint sensitivity across penalty values (RBF model).}
\label{tab:mst_pen}
\begin{tabular}{lcc}
\toprule
Penalty & Reddit breakpoints & Hacker News breakpoints \\
\midrule
0.5 & T4; T6 & T3; T6 \\
1.0 & T4; T6 & T3; T6 \\
2.0 & T4     & T4 \\
3.0 & none   & none \\
5.0 & none   & none \\
10.0 & none  & none \\
\bottomrule
\end{tabular}
\end{table}

\begin{figure}[H]
\centering
\includegraphics[width=0.85\linewidth]{mst_penalty_sensitivity.png}
\caption{PELT changepoint sensitivity across penalties 0.5--10. Vertical dotted lines
mark detected breakpoints per penalty; no single breakpoint location is robust across
settings.}
\end{figure}

We distinguish two \emph{heuristic} operationalizations suggested by the curves
themselves (independent of any single penalty choice):

\begin{itemize}
  \item \textbf{Changepoint MST (penalty-dependent)}: often T4 (50 comments) at
        penalty~$= 2$, but not robust.
  \item \textbf{Practical MST}: T5--T6 (100--200 comments) --- where accuracy
        effectively plateaus and marginal gains fall below 0.06\% per comment.
\end{itemize}

On Reddit, 50 comments corresponds to approximately 2{,}000--3{,}000 words and
(given the average posting rate in our sample) roughly 2--4 months of typical
activity. This level marks a region of accelerating gains in our curves, not a
hard threshold of deanonymizability.

At the Practical MST of 100--200 comments (T5--T6), this corresponds to roughly
4--8 months of typical Reddit activity. Beyond this region, marginal accuracy
gains become very small despite substantial increases in available text.


\section{Calibration Analysis}


A striking finding emerges from the calibration analysis. Despite top-1 accuracy
reaching nearly 40\%, \textbf{Recall@90\% Precision remains extremely low}:

\begin{itemize}
  \item Reddit: Recall@90\% Precision $\leq$ 2.1\% across all truncation levels
  \item Hacker News: Recall@90\% Precision $\leq$ 0.62\% across all truncation levels
\end{itemize}

This near-zero recall at high precision is explained by severe overconfidence in
GPT-4o's confidence estimates:

\begin{table}[H]
\centering
\caption{GPT-4o confidence vs.\ empirical accuracy on Hacker News.}
\begin{tabular}{cc}
\toprule
Stated Confidence & Empirical Accuracy (\%) \\
\midrule
0.80 & 14.4 \\
0.85 & 26.7 \\
0.90 & 49.0 \\
\bottomrule
\end{tabular}
\end{table}

At a stated confidence of 0.85, the model is empirically correct only 26.7\% of the
time---a gap of over 58 percentage points. Confidence scores cluster tightly in the
0.80--0.90 range, providing almost no discrimination between correct and incorrect
predictions. As a result, no confidence threshold can simultaneously achieve high
precision and non-trivial recall.

\begin{figure}[H]
\centering
\includegraphics[width=0.85\linewidth]{confidence_calibration_hn.png}
\caption{GPT-4o confidence vs.\ empirical accuracy on Hacker News. The model is
systematically overconfident, with stated confidence far exceeding actual accuracy
across all confidence levels.}
\end{figure}

This finding has practical implications: an automated ESRC-based attacker that
filters on confidence would discard the vast majority of correct predictions. The
attack is substantially less dangerous in operational settings than top-1 accuracy
alone suggests.


\section{Search vs.\ Reason Analysis}


To understand which stage of the pipeline limits performance, we compare retrieval recall (the fraction of true matches present in the top-15 candidates returned by FAISS) against final top-1 accuracy. We additionally report Conditional Reason Accuracy, defined as top-1 accuracy divided by Recall@15—that is, the fraction of queries for which the true candidate was retrieved and subsequently selected correctly by GPT-4o.

At T8, Conditional Reason Accuracy reaches 73.2\% for Reddit and 64.0\% for Hacker News.

\begin{table}[H]
\centering
\caption{Search-stage vs.\ Reason-stage performance at T8.}
\begin{tabular}{lccc}
\toprule
Dataset & Recall@15 (Search) & Top-1 Ac. (Reason)  \\
\midrule
Reddit       & 53.8\% & 39.4\% \\
Hacker News  & 59.4\% & 38.0\% \\
\bottomrule
\end{tabular}
\end{table}

At T8, FAISS retrieves the true match within the top-15 for 53.8\% of Reddit queries
and 59.4\% of Hacker News queries. GPT-4o then selects the correct user in 73.2\% and
64.0\% of those retrievable cases, respectively---but in only 39.4\% and 38.0\% of
all queries overall. The Reason stage therefore fails on a substantial fraction of
queries even when search succeeds; improving retrieval alone would not close the full
gap between Recall@15 and top-1 accuracy.

\begin{figure}[H]
\centering
\includegraphics[width=0.85\linewidth]{search_vs_reason_hn.png}
\caption{Comparison of Search-stage retrieval recall and Reason-stage top-1 accuracy
across truncation levels on Hacker News.}
\end{figure}

This result identifies the \textbf{Reason stage as the primary bottleneck} of the
ESRC pipeline. Improving embedding quality or expanding the candidate set would yield
limited gains relative to improving the reasoning and selection step---for example
through chain-of-thought prompting, finer-grained profile representations, or
re-ranking strategies.

\subsection{Diversity Experiment}

\begin{table}[H]
\centering
\caption{Diversity experiment: top-1 accuracy and FAISS Recall@15 by subreddit
diversity group (T4, 50 comments).}
\label{tab:diversity}
\begin{tabular}{lcccc}
\toprule
Group & $n$ & Top-1  & 95\% CI  & R@15  \\
\midrule
Low (1--3)    & 58  & 27.59 & [17.24, 39.66] & 51.7 \\
Medium (4--10) & 179 & 21.23 & [15.08, 27.37] & 43.0 \\
High (11+)    & 263 & 19.01 & [14.45, 23.95] & 35.4 \\
\bottomrule
\end{tabular}
\end{table}

Top-1 accuracy decreases from 27.59\% (low) to 19.01\% (high). The pattern is
already visible at the Search stage: FAISS Recall@15 falls from 51.7\% to 35.4\%.
Pairwise two-proportion $z$-tests are not significant at $\alpha = 0.05$ (low vs.\
medium: $p = 0.316$; low vs.\ high: $p = 0.144$; medium vs.\ high: $p = 0.567$),
and bootstrap 95\% confidence intervals overlap substantially
(Table~\ref{tab:diversity}).


\begin{figure}[H]
\centering
\includegraphics[width=0.85\linewidth]{diversity_top1_accuracy_with_ci.png}
\caption{Diversity experiment: top-1 accuracy with bootstrap 95\% confidence
intervals (10{,}000 resamples). Overlapping intervals indicate that group
differences are not statistically separable at this sample size.}
\end{figure}

\subsection{Content-Type Experiment}

Of 500 classified users, 447 qualify for $\geq 50$ Personal, 499 for Opinion, 500
for Topical, and 447 for all three types. The ESRC analysis uses the 432 users with
complete Extract summaries for all three types.

\begin{table}[H]
\centering
\caption{Content-type experiment: top-1 accuracy, bootstrap 95\% confidence intervals, and FAISS Recall@15 (ESRC cohort with all three summaries).}
\label{tab:content_type}
\begin{tabular}{lccccc}
\toprule
Type & $n$ & Top-1 & 95\% CI & R@15 \\
\midrule
Personal & 432 & 32.41 & [28.01, 36.81] & 46.3 \\
Opinion  & 426  & 18.78 & [15.26, 22.54] & 33.1 \\
Topical  & 432  & 18.29 & [14.81, 21.99] & 35.2 \\
\bottomrule
\end{tabular}
\end{table}

Personal-only profiles substantially outperform Opinion-only and Topical-only
profiles. Pairwise tests are significant for Personal vs.\ Opinion ($p = 0.000005$)
and Personal vs.\ Topical ($p = 0.000002$); Opinion vs.\ Topical is not significant
($p = 0.853$). The effect appears already in Search: Personal Recall@15 is 46.3\%,
compared with 33.1\% (Opinion) and 35.2\% (Topical). Only Personal content achieves
non-trivial Recall@90\% Precision (1.43\%); no confidence threshold for Opinion or
Topical reaches 90\% precision.

Personal self-disclosure carries disproportionate identifying information. At the
same fixed volume of 50 comments, Personal-only profiles achieve 32.4\% top-1
accuracy, compared with approximately 18\% for Opinion-only and Topical-only
profiles.

\begin{figure}[H]
\centering
\includegraphics[width=0.85\linewidth]{content_type_top1_accuracy_with_ci.png}
\caption{Content-type experiment: top-1 accuracy with bootstrap 95\% confidence
intervals. Personal content is significantly more deanonymizable than Opinion or
Topical content ($p < 0.00001$).}
\end{figure}


\section{Discussion}


\paragraph{On the practical significance of the MST.}
Our exploratory changepoint analysis can suggest an acceleration region near 50
comments under specific penalty settings, with practical saturation near 100--200
comments where marginal gains fall below 0.06\% per comment. Because breakpoint
locations vary with the penalty parameter, we treat 50 comments as one plausible
interpretation of the curve shape rather than a validated threshold. In practical
terms, users with more than a few months of typical activity are likely well above
the region where accuracy gains are largest.

\paragraph{On the gap between accuracy and precision-constrained recall.}
The near-zero Recall@90\% Precision values present an apparent paradox: how can an
attack reach 39\% accuracy yet remain nearly unusable at high precision? The answer
lies in the structure of confidence scores. GPT-4o assigns narrow, near-identical
confidence values to correct and incorrect predictions alike. An attacker who accepts
only high-confidence predictions is left with almost nothing. This does not make the
attack safe---a 39\% unconditional success rate is substantial---but it does suggest
that confidence-gated deployment of such attacks is severely limited by current model
calibration.

\paragraph{On cross-platform generalizability.}
Strong Spearman correlations (Reddit $\rho = 0.976$; Hacker News $\rho = 1.000$;
both $p < 0.001$ over truncation level ranks) and a convergent accuracy ceiling
($\sim$38--40\%) on two very different platforms suggest that the
relationship between text volume and deanonymizability is largely a property of the
ESRC method rather than any single platform's characteristics, although bootstrap
intervals indicate that platform gaps at individual truncation levels are often
within sampling noise.

\paragraph{On community diversity.}
Grouping users by unique subreddits at fixed T4 volume yields a downward trend in
both FAISS Recall@15 and top-1 accuracy from low- to high-diversity users. One
plausible explanation is that narrow community focus produces a clearer topical
identity signal, whereas users spread across many subreddits present a more diffuse
profile. However, confidence intervals overlap and all pairwise significance tests
fail to reject the null at $\alpha = 0.05$. We therefore treat diversity as a
suggestive but inconclusive moderator---not evidence that community breadth
significantly affects deanonymization risk.

\paragraph{On content type.}
The content-type experiment provides the strongest additional result beyond
truncation volume. Personal disclosures are significantly more identifying than
opinion statements or topical discussion at the same 50-comment volume, with the
advantage emerging already at retrieval. This implies that not all text carries equal
privacy risk: a user who primarily discusses topics may be less exposed than one who
routinely shares personal facts, even at identical comment counts. Platform guidance
and user education should distinguish \textit{what} is shared from \textit{how much}
is shared.

\paragraph{Limitations.}
Our evaluation uses a closed candidate pool of 1{,}000 users. In open-world settings,
accuracy would be substantially lower. Additionally, all experiments use the Lermen
G.2 prompts without modification; prompt optimization could shift results in either
direction. The use of GPT-4o for the Reason stage introduces a dependency on a
proprietary model whose behavior may change across versions.


\section{Conclusion}


We present a systematic study of the relationship between text volume and LLM-based
deanonymization performance. Across Reddit and Hacker News, top-1 accuracy increases
nearly monotonically with text volume (Reddit $\rho = 0.976$; Hacker News
$\rho = 1.000$; both $p < 0.001$ over truncation level ranks), reaching
approximately 39--40\% at full profile size, with bootstrap confidence intervals
indicating that adjacent-level differences should be interpreted cautiously.

Exploratory changepoint analysis suggests a possible acceleration near 50 comments
under penalty~$= 2$, but breakpoint locations are penalty-sensitive and should not be
reported as a definitive MST. Practical saturation appears between 100 and 200
comments, where marginal gains fall below 0.06\% per comment.

Three findings extend beyond the core truncation result. First, GPT-4o is severely
miscalibrated: at stated confidence 0.85, it is correct only 26.7\% of the time,
limiting the precision-constrained utility of the attack. Second, the Reason stage is
the primary bottleneck: conditional Reason accuracy reaches only 64--73\% even when
FAISS retrieves the true match in the top-15. Third, the accuracy ceiling of
$\sim$40\% is consistent across two platforms with substantially different writing
cultures, suggesting a method-level constraint rather than a platform-specific one.

Fourth, at fixed T4 volume, \textbf{content type matters}: Personal-only profiles
reach 32.4\% top-1 accuracy versus $\sim$18\% for Opinion-only and Topical-only
profiles ($p < 0.00001$), indicating that personal self-disclosure carries
disproportionate identifying signal. Fifth, a community-diversity grouping shows a statistically
non-significant inverse trend---users in fewer subreddits appear somewhat more
identifiable---but overlapping confidence intervals preclude a definitive claim.

Together, these findings show not only \textit{how much} text matters for
LLM-based deanonymization, but also \textit{what type} of text matters. Personal
disclosure is far more identifying than topical participation at equal volume.

These findings contribute to a quantitative understanding of textual privacy risk,
and point toward calibration improvement and reasoning-stage optimization as the
most promising directions for both strengthening and defending against LLM-based
deanonymization.


\bibliographystyle{ACM-Reference-Format}
\bibliography{references}

\end{document}